\begingroup
\shorthandoff{"<>"}
\setlength{\arrayrulewidth}{1pt}
\arrayrulecolor{vlepsbased}
\begin{longtable}{|p{15cm}|}
  \hline
  \rowcolor{maincolor} 
  \multicolumn{1}{|l|}{\textcolor{white}{\textbf{Caso de Uso}}} \\
  \hline
  \endfirsthead
  \multicolumn{1}{c}{{\tablename\ \thetable{} -- Continuación}} \\
  \hline
  \rowcolor{maincolor} 
  \multicolumn{1}{|l|}{\textcolor{white}{\textbf{Caso de Uso (Continuación)}}} \\
  \hline
  \endhead
  \hline \multicolumn{1}{|r|}{\textit{Continúa en la siguiente página...}} \\ \hline
  \endfoot
  \hline
  \caption{Caso de Uso CU10 — Renombrado manual de sesión de chat.}
  \label{TB:CU10_RENOMBRAR_SESION}
  \endlastfoot

  \rowcolor{ulepsbased}
  \begin{tabularx}{\linewidth}{@{}X X@{}}
    \textbf{Identificador:} CU10 & \textbf{Actor principal:} Usuario
  \end{tabularx} \\ \hline

  \textbf{Nombre:} Renombrado manual de sesión de chat \\ \hline

  \rowcolor{ulepsbased}
  \textbf{Descripción:} \newline 
  El usuario decide modificar el título de una sesión de chat existente en su panel lateral para facilitar su organización, introduciendo un nuevo nombre. \\ \hline

  \textbf{Precondiciones:} \newline
  1. El usuario ha iniciado sesión. \newline
  2. Existe al menos una sesión de chat en su historial. \\ \hline

  \rowcolor{ulepsbased}
  \textbf{Flujo principal:} \newline
  1. El usuario hace clic en el botón de "Renombrar" junto a la sesión deseada en el panel lateral. \newline
  2. La interfaz reemplaza el título estático por un campo de texto editable y solicita el nuevo título. \newline
  3. El usuario escribe el nuevo nombre y confirma (pulsando Enter o el botón de guardado). \newline
  4. La interfaz valida que el campo de texto no esté vacío ni compuesto solo por espacios en blanco. \newline
  5. La interfaz envía la petición al backend con el nuevo título. \newline
  6. El backend verifica la propiedad de la sesión y comprueba que el título no llegue vacío. \newline
  7. Si el título supera los 200 caracteres, el backend realiza un truncado silencioso recortándolo al límite permitido. \newline
  8. El sistema actualiza el título en la base de datos y devuelve el título modificado exitosamente. \newline
  9. El frontend actualiza el nombre visualmente en el panel lateral. \\ \hline

  \textbf{Flujos alternativos:} \newline
  \textbf{FA-01:} Título vacío o en blanco: \newline
  1. El sistema detecta en el frontend que el título introducido está vacío o contiene únicamente espacios en blanco. \newline
  2. La función de guardado se aborta sin llegar a enviar ninguna solicitud HTTP al backend. \newline
  3. La interfaz cancela el modo de edición silenciosamente y restaura el título original. \\ \hline

  \rowcolor{ulepsbased}
  \textbf{Postcondiciones:} \newline
  1. El nuevo título queda asociado a la sesión en la base de datos de manera permanente. \\ \hline

\end{longtable}
\endgroup
